% ============================================================
%  طرح پیشنهادی مجلس مهستان — نسخه‌ی خلاصه‌ی اجرایی
%  کامپایل با: XeLaTeX
%  فونت مورد نیاز: IRNazanin یا هر فونت فارسی دلخواه
% ============================================================
\documentclass[a4paper,11pt]{article}

% ---------- بسته‌های پایه ----------
\usepackage[top=1.8cm, bottom=2cm, left=2cm, right=2cm]{geometry}
\usepackage{fontspec}
\usepackage{graphicx}
\usepackage{xcolor}
\usepackage{tikz}
\usetikzlibrary{shapes.geometric, arrows.meta, positioning, calc, decorations.pathmorphing, backgrounds, fit, shadows.blur}
\usepackage{tcolorbox}
\tcbuselibrary{skins, breakable, hooks}
\usepackage{tabularx}
\usepackage{booktabs}
\usepackage{multicol}
\usepackage{enumitem}
\usepackage{fancyhdr}
\usepackage{hyperref}
\usepackage{pifont}
\usepackage{amssymb}
\usepackage{colortbl}
\usepackage{array}
\usepackage{ragged2e}
\usepackage{setspace}
\usepackage{titlesec}
\usepackage{fancybox}
\usepackage{mdframed}
\usepackage{lastpage}

% ---------- تعریف رنگ‌ها ----------
\definecolor{mahestan-blue}{HTML}{1565C0}
\definecolor{mahestan-green}{HTML}{2E7D32}
\definecolor{mahestan-gold}{HTML}{F9A825}
\definecolor{mahestan-red}{HTML}{C62828}
\definecolor{mahestan-purple}{HTML}{6A1B9A}
\definecolor{mahestan-orange}{HTML}{E65100}
\definecolor{mahestan-teal}{HTML}{00838F}
\definecolor{bg-light-blue}{HTML}{E3F2FD}
\definecolor{bg-light-green}{HTML}{E8F5E9}
\definecolor{bg-light-gold}{HTML}{FFF9C4}
\definecolor{bg-light-red}{HTML}{FFEBEE}
\definecolor{bg-light-purple}{HTML}{F3E5F5}
\definecolor{bg-cream}{HTML}{FFFDE7}
\definecolor{bg-gray}{HTML}{F5F5F5}
\definecolor{border-gray}{HTML}{BDBDBD}
\definecolor{text-dark}{HTML}{212121}
\definecolor{text-medium}{HTML}{424242}

% ---------- تنظیم هایپرلینک ----------
\hypersetup{
    colorlinks=true,
    linkcolor=mahestan-blue,
    urlcolor=mahestan-teal,
    pdfTitle={طرح پیشنهادی مجلس مهستان},
    pdfAuthor={},
    pdfsubject={چارچوب نهادی گذار دموکراتیک ایران},
}

% ---------- تعریف باکس‌ها ----------
\newtcolorbox{principlebox}[1][]{
    enhanced, breakable,
    colback=bg-light-blue,
    colframe=mahestan-blue,
    coltitle=white,
    fonttitle=\bfseries\large,
    boxrule=1.5pt,
    arc=4pt,
    left=8pt, right=8pt, top=6pt, bottom=6pt,
    shadow={2pt}{-2pt}{0pt}{black!20},
    #1
}

\newtcolorbox{alertbox}[1][]{
    enhanced, breakable,
    colback=bg-light-red,
    colframe=mahestan-red,
    coltitle=white,
    fonttitle=\bfseries,
    boxrule=1.5pt,
    arc=4pt,
    left=8pt, right=8pt, top=6pt, bottom=6pt,
    #1
}

\newtcolorbox{quotebox}[1][]{
    enhanced, breakable,
    colback=bg-cream,
    colframe=mahestan-gold,
    boxrule=0pt,
    borderline west={3pt}{0pt}{mahestan-gold},
    left=12pt, right=8pt, top=6pt, bottom=6pt,
    sharp corners,
    #1
}

\newtcolorbox{greenbox}[1][]{
    enhanced, breakable,
    colback=bg-light-green,
    colframe=mahestan-green,
    coltitle=white,
    fonttitle=\bfseries,
    boxrule=1.5pt,
    arc=4pt,
    left=8pt, right=8pt, top=6pt, bottom=6pt,
    #1
}

\newtcolorbox{purplebox}[1][]{
    enhanced, breakable,
    colback=bg-light-purple,
    colframe=mahestan-purple,
    coltitle=white,
    fonttitle=\bfseries,
    boxrule=1.5pt,
    arc=4pt,
    left=8pt, right=8pt, top=6pt, bottom=6pt,
    #1
}

% ---------- تنظیمات عنوان‌ها ----------
\titleformat{\section}
    {\color{mahestan-blue}\bfseries\Large}
    {\thesection.}{0.5em}{}
    [\vspace{-0.5em}\textcolor{mahestan-blue}{\rule{\textwidth}{1.5pt}}]

\titleformat{\subsection}
    {\color{mahestan-purple}\bfseries\large}
    {\thesubsection.}{0.5em}{}

\titleformat{\subsubsection}
    {\color{mahestan-green}\bfseries\normalsize}
    {\thesubsubsection.}{0.5em}{}

% ---------- سربرگ و پابرگ ----------
\pagestyle{fancy}
\fancyhf{}
\renewcommand{\headrulewidth}{1pt}
\renewcommand{\headrule}{\hbox to\headwidth{\color{mahestan-blue}\leaders\hrule height \headrulewidth\hfill}}
\renewcommand{\footrulewidth}{0.5pt}
\renewcommand{\footrule}{\hbox to\headwidth{\color{border-gray}\leaders\hrule height \footrulewidth\hfill}}

% ---------- بسته‌ی فارسی (باید آخرین بسته باشد) ----------
\usepackage{xepersian}
\settextfont[Scale=1.05]{vazirmatn}
\setdigitfont{vazirmatn}
\DefaultMathsDigits

\fancyhead[R]{\small\textcolor{mahestan-blue}{\textbf{طرح پیشنهادی مجلس مهستان}}}
\fancyhead[L]{\small\textcolor{text-medium}{نسخه‌ی خلاصه‌ی اجرایی}}
\fancyfoot[C]{\small\textcolor{text-medium}{صفحه‌ی \thepage\ از \pageref{LastPage}}}
\fancyfoot[L]{\tiny\textcolor{border-gray}{اسفند ۱۴۰۴}}

% ---------- شروع سند ----------
\begin{document}
\setstretch{1.3}

% ==================== صفحه‌ی عنوان ====================
\thispagestyle{empty}

\begin{tikzpicture}[remember picture, overlay]
    % پس‌زمینه‌ی بالا
    \fill[mahestan-blue] (current page.north west) rectangle ([yshift=-6cm]current page.north east);
    % نوار طلایی
    \fill[mahestan-gold] ([yshift=-6cm]current page.north west) rectangle ([yshift=-6.4cm]current page.north east);
    % نوار سبز
    \fill[mahestan-green!70] ([yshift=-6.4cm]current page.north west) rectangle ([yshift=-6.6cm]current page.north east);
    % دایره‌ی تزئینی
    \node[circle, fill=white, minimum size=2.5cm, inner sep=0pt,
          blur shadow={shadow blur steps=15, shadow xshift=0pt, shadow yshift=-2pt}]
          at ([yshift=-6cm]current page.north) {\textcolor{mahestan-blue}{\Huge\textbf{🏛}}};
    % خطوط تزئینی پایین
    \fill[mahestan-blue!10] (current page.south west) rectangle ([yshift=2.5cm]current page.south east);
    \fill[mahestan-gold] ([yshift=2.5cm]current page.south west) rectangle ([yshift=2.7cm]current page.south east);
\end{tikzpicture}

\vspace*{1.5cm}

\begin{center}
    {\Huge\textcolor{white}{\textbf{طرح پیشنهادی}}}\\[0.3cm]
    {\fontsize{36}{42}\selectfont\textcolor{white}{\textbf{مجلس مهستان}}}\\[0.5cm]
    {\Large\textcolor{bg-light-blue}{چارچوب نهادی گذار دموکراتیک ایران}}
\end{center}

\vspace{3cm}

\begin{center}
    \begin{tcolorbox}[
        enhanced,
        width=12cm,
        colback=bg-cream,
        colframe=mahestan-gold,
        boxrule=2pt,
        arc=8pt,
        halign=center,
        shadow={3pt}{-3pt}{0pt}{black!15},
    ]
    {\large\textbf{نسخه‌ی خلاصه‌ی اجرایی}}\\[0.3cm]
    {\normalsize ایران برای همه‌ی ایرانیان}\\[0.3cm]
    {\small\textcolor{text-medium}{اسفند ۱۴۰۴ / فوریه‌ی ۲۰۲۶}}
    \end{tcolorbox}
\end{center}

\vspace{2cm}

\begin{center}
    \textcolor{mahestan-purple}{\rule{8cm}{0.5pt}}\\[0.5cm]
    {\small\textcolor{text-medium}{%
    این سند زنده است و با مشارکت همگان بهتر خواهد شد.\\
    تمام اعداد و نسبت‌ها پیشنهادی هستند و توسط مجلس منتخب مردم نهایی خواهند شد.}}
\end{center}

\newpage

% ==================== فهرست مطالب ====================
\thispagestyle{fancy}

\begin{purplebox}[title={\large فهرست مطالب}]
\begin{tabularx}{\textwidth}{>{\centering\bfseries}p{1cm} X >{\centering}p{1.5cm}}
\textcolor{mahestan-purple}{بخش} & \textcolor{mahestan-purple}{\bfseries موضوع} & \textcolor{mahestan-purple}{\bfseries صفحه} \\
\midrule
۱ & چرایی و ضرورت طرح مهستان & ۲ \\
۲ & شش اصل راهنمای طراحی & ۲ \\
۳ & ساختار مجلس مهستان و ترکیب آن & ۳ \\
۴ & نهادهای کلیدی دوران گذار و عدالت انتقالی & ۳ \\
۵ & اولویت‌های ۱۰۰ روز نخست & ۴ \\
۶ & اصول بنیادین غیرقابل‌نقض & ۴ \\
۷ & نقشه‌ی راه زمانی & ۵ \\
۸ & فراخوان مشارکت & ۵ \\
\end{tabularx}
\end{purplebox}

\vspace{0.8cm}

% ==================== بخش ۱ ====================
\section{چرایی و ضرورت}

\begin{quotebox}
{\itshape «کیفیت طراحی نهادی دوران گذار، بیش از هر عامل دیگری، سرنوشت نظام سیاسی آینده را تعیین می‌کند.»}
\end{quotebox}

\vspace{0.3cm}

ایران در مقطعی تاریخی قرار دارد. تجربه‌ی جهانی یک درس بنیادین دارد:

\begin{center}
\begin{tikzpicture}
    \node[draw=mahestan-green, fill=bg-light-green, rounded corners=6pt,
          minimum width=7cm, minimum height=1.2cm, text width=6.5cm, align=center,
          font=\small\bfseries]
        (success) {✅ گذار طراحی‌شده + فراگیر $\longrightarrow$ دموکراسی پایدار\\
        {\footnotesize\textcolor{text-medium}{آفریقای جنوبی · اسپانیا · شیلی · لهستان}}};
    \node[draw=mahestan-red, fill=bg-light-red, rounded corners=6pt,
          minimum width=7cm, minimum height=1.2cm, text width=6.5cm, align=center,
          font=\small\bfseries, below=0.4cm of success]
        (fail) {❌ گذار شتاب‌زده + حذفی $\longrightarrow$ هرج‌ومرج یا استبداد نو\\
        {\footnotesize\textcolor{text-medium}{لیبی · مصر · عراق}}};
\end{tikzpicture}
\end{center}

مردم ایران — از هر قشر و طبقه — هزینه‌های گزافی پرداخت کرده‌اند. آینده باید با \textbf{اندیشه، اجماع و مهندسی نهادی} ساخته شود؛ نه با تعلل و ترس از تغییر، و نه با شتاب و تمامیت‌خواهی.

% ==================== بخش ۲ ====================
\section{شش اصل راهنمای طراحی}

\begin{center}
\begin{tikzpicture}[
    principle/.style={
        draw=#1, fill=#1!8, rounded corners=5pt,
        minimum width=4.8cm, minimum height=1.7cm,
        text width=4.4cm, align=center, font=\small,
        blur shadow={shadow blur steps=5, shadow xshift=1pt, shadow yshift=-1pt}
    }
]
    \node[principle=mahestan-blue]  (p1) at (0,0)
        {\textbf{\textcolor{mahestan-blue}{اصل ۱}}\\ترکیب دموکراسی\\مستقیم و تخصص‌گرایی};
    \node[principle=mahestan-green] (p2) at (5.5,0)
        {\textbf{\textcolor{mahestan-green}{اصل ۲}}\\نمایندگی نسبتی\\هیچ صدایی نشنیده نماند};
    \node[principle=mahestan-purple](p3) at (11,0)
        {\textbf{\textcolor{mahestan-purple}{اصل ۳}}\\قانون‌اساسی‌گرایی\\دو مرحله‌ای (موقت + نهایی)};
    \node[principle=mahestan-orange](p4) at (0,-2.5)
        {\textbf{\textcolor{mahestan-orange}{اصل ۴}}\\اصلاح نه تخریب\\حفظ نهادها با رهبری جدید};
    \node[principle=mahestan-teal]  (p5) at (5.5,-2.5)
        {\textbf{\textcolor{mahestan-teal}{اصل ۵}}\\شفافیت و پاسخ‌گویی\\حداکثری};
    \node[principle=mahestan-red]   (p6) at (11,-2.5)
        {\textbf{\textcolor{mahestan-red}{اصل ۶}}\\فراگیری و تعهد\\به حقوق بنیادین};
\end{tikzpicture}
\end{center}

\vspace{0.3cm}

\subsection{اصل محوری: تفکیک واقعیت از ارزش}

\begin{center}
\begin{tikzpicture}[
    box/.style={draw=#1, fill=#1!8, rounded corners=4pt,
        minimum width=5cm, minimum height=2.2cm, text width=4.6cm,
        align=center, font=\small}
]
    \node[box=mahestan-orange] (fact) at (0,0)
        {\textbf{\large 🔬 بُعد واقعیت}\\[2pt]
        \textcolor{text-medium}{وضعیت چیست؟}\\[2pt]
        منبع: متخصصان و پژوهشگران\\
        ابزار: مرکز تحقیقات مجلس};
    \node[box=mahestan-blue] (value) at (7,0)
        {\textbf{\large 📢 بُعد ارزش}\\[2pt]
        \textcolor{text-medium}{مردم چه می‌خواهند؟}\\[2pt]
        منبع: آرای مستقیم مردم\\
        ابزار: واحد ارتباط مردمی};
    \node[draw=mahestan-green, fill=bg-light-green, rounded corners=4pt,
        minimum width=3.5cm, minimum height=1.3cm, align=center, font=\small\bfseries]
        (decision) at (3.5,-2.5)
        {🏛️ تصمیم‌گیری آگاهانه\\واقعیت + ارزش = حکمرانی مطلوب};
    \draw[-{Stealth[length=5pt]}, thick, mahestan-orange] (fact.south) -- (decision.north west);
    \draw[-{Stealth[length=5pt]}, thick, mahestan-blue]  (value.south) -- (decision.north east);
\end{tikzpicture}
\end{center}

{\footnotesize\textcolor{text-medium}{%
\textbf{مبانی نظری:} هیوم (تفکیک هست/باید) · هابرماس (دموکراسی مشورتی) · لیپهارت (الگوهای دموکراسی)
\quad\textbf{الگوهای عملی:} مجمع شهروندان ایرلند · کنوانسیون شهروندی فرانسه · فرایند قانون اساسی آفریقای جنوبی}}

% ==================== بخش ۳ ====================
\section{ساختار مجلس مهستان}

\subsection{ترکیب: ۳۰۰ نماینده‌ی منتخب مردم}

\begin{center}
\begin{tikzpicture}
    % Background
    \fill[bg-light-blue, rounded corners=8pt] (-0.5,-0.3) rectangle (13,3.2);
    \node[font=\bfseries\large, anchor=north] at (6.25,3) {🗳️ ۳۰۰ کرسی مجلس مهستان};
    % Bars
    \fill[mahestan-blue, rounded corners=2pt] (0.5,1.6) rectangle (10.5,2.2);
    \node[font=\small\bfseries, white] at (5.5,1.9) {۲۴۰ کرسی استانی — نسبی-فهرستی — آستانه ۳٪};
    \fill[mahestan-green, rounded corners=2pt] (0.5,0.8) rectangle (5,1.4);
    \node[font=\small\bfseries, white] at (2.75,1.1) {۴۰ کرسی فهرست ملی — آستانه ۵٪};
    \fill[mahestan-purple, rounded corners=2pt] (0.5,0) rectangle (3,0.6);
    \node[font=\small\bfseries, white] at (1.75,0.3) {۲۰ کرسی تضمینی اقلیت‌ها};
    % Notes
    \node[font=\footnotesize, text=text-medium, anchor=west] at (5.3,1.1) {سراسری};
    \node[font=\footnotesize, text=text-medium, anchor=west] at (3.3,0.3) {قومی · زبانی · مذهبی};
\end{tikzpicture}
\end{center}

\vspace{0.1cm}

\begin{center}
{\small
\textcolor{mahestan-red}{⚖️ حداقل ۳۰٪ تنوع جنسیتی الزامی (قاعده‌ی ۱ به ۳)}
\quad\textbullet\quad
\textcolor{mahestan-teal}{🌐 سازوکار مشارکت ایرانیان خارج از کشور}
}
\end{center}

\subsection{تشکیل کابینه‌ی اجرایی موقت}

\begin{center}
\begin{tabularx}{0.85\textwidth}{>{\centering\bfseries\small}p{3.5cm} >{\small\RaggedRight}X}
\toprule
\rowcolor{bg-light-blue}
\textcolor{mahestan-blue}{سناریو} & \textcolor{mahestan-blue}{\bfseries شیوه‌ی تشکیل کابینه} \\
\midrule
اکثریت مطلق (۵۰٪+۱) & آن جریان مسئول تشکیل کابینه خواهد بود \\
\rowcolor{bg-gray}
بدون اکثریت مطلق & ائتلاف جریان‌ها تا رسیدن به اکثریت قاطع \\
عدم ائتلاف & رأی‌گیری مجلس برای هر پست وزارتی به‌طور جداگانه \\
\bottomrule
\end{tabularx}
\end{center}

{\footnotesize\textcolor{text-medium}{%
\textbf{پاسخ‌گویی کابینه:}
رأی اعتماد اولیه: ۵۰٪+۱
\quad|\quad
استیضاح و برکناری: ۷۵٪
\quad|\quad
تعویض فردی وزرا: ⅔ آرا
}}

\subsection{یازده کمیسیون تخصصی دائمی}

\begin{center}
\begin{tikzpicture}[
    comm/.style={draw=border-gray, fill=bg-gray, rounded corners=3pt,
        minimum width=3.3cm, minimum height=0.7cm, align=center, font=\tiny\bfseries}
]
    \node[comm, fill=mahestan-blue!10]  at (0,1.2)   {📜 قانون اساسی موقت};
    \node[comm, fill=mahestan-red!10]   at (3.8,1.2)  {⚖️ عدالت انتقالی و حقوق بشر};
    \node[comm, fill=mahestan-orange!10]at (7.6,1.2)  {🛡️ اصلاح امنیتی و نظامی};
    \node[comm, fill=mahestan-purple!10]at (11.4,1.2) {⚔️ اصلاح قضایی و حقوقی};
    \node[comm, fill=mahestan-green!10] at (0,0.3)    {💰 اقتصاد و بودجه};
    \node[comm, fill=mahestan-teal!10]  at (3.8,0.3)  {📚 آموزش، علم و فرهنگ};
    \node[comm, fill=mahestan-blue!10]  at (7.6,0.3)  {🌐 سیاست خارجی};
    \node[comm, fill=mahestan-gold!15]  at (11.4,0.3) {🏘️ حقوق اقلیت‌ها};
    \node[comm, fill=mahestan-green!10] at (1.9,-0.6) {🌿 محیط زیست};
    \node[comm, fill=mahestan-red!10]   at (5.7,-0.6) {👩 حقوق زنان و برابری};
    \node[comm, fill=mahestan-purple!10]at (9.5,-0.6) {📺 رسانه و حقوق دیجیتال};
\end{tikzpicture}
\end{center}

% ==================== بخش ۴ ====================
\section{نهادهای کلیدی دوران گذار و عدالت انتقالی}

\begin{center}
\begin{tikzpicture}[
    node distance=0.6cm,
    inst/.style={draw=#1, fill=#1!8, rounded corners=4pt,
        minimum width=4cm, minimum height=0.8cm, align=center, font=\small}
]
    % مرکز
    \node[draw=mahestan-green, fill=bg-light-green, rounded corners=6pt,
        minimum width=6cm, minimum height=1.2cm, align=center,
        font=\bfseries, blur shadow={shadow blur steps=5}]
        (meh) {🏛️ مجلس مهستان\\[-2pt]{\footnotesize نهاد عالی حاکمیت دوران گذار}};

    % ردیف اول
    \node[inst=mahestan-gold, below left=0.8cm and 0.5cm of meh]  (cab)
        {⚙️ کابینه‌ی اجرایی موقت};
    \node[inst=mahestan-red, below=0.8cm of meh]  (jud)
        {⚖️ شورای قضایی انتقالی};
    \node[inst=mahestan-orange, below right=0.8cm and 0.5cm of meh]  (mil)
        {🛡️ فرماندهی نظامی-انتظامی};

    % ردیف دوم
    \node[inst=mahestan-purple, below=0.5cm of cab]  (ssr) {🔄 کمیسیون اصلاح امنیتی};
    \node[inst=mahestan-blue, below=0.5cm of jud]  (trc) {📋 کمیسیون حقیقت‌یابی};
    \node[inst=mahestan-teal, below=0.5cm of mil]  (elec) {🗳️ هیأت نظارت انتخابات};

    % ردیف سوم
    \node[inst=mahestan-green, below=0.5cm of ssr]  (omb) {👤 آمبودزمان};
    \node[inst=mahestan-blue, below=0.5cm of trc]  (med) {📺 کمیسیون ملی رسانه‌ها};
    \node[inst=mahestan-orange, below=0.5cm of elec]  (econ) {📊 شورای اقتصادی};

    % فلش‌ها
    \draw[-{Stealth}, thick, mahestan-green] (meh) -- (cab);
    \draw[-{Stealth}, thick, mahestan-green] (meh) -- (jud);
    \draw[-{Stealth}, thick, mahestan-green] (meh) -- (mil);
\end{tikzpicture}
\end{center}

\subsection{عدالت انتقالی — پنج ستون}

\begin{center}
\begin{tikzpicture}[
    pillar/.style={draw=#1, fill=#1!10, rounded corners=4pt,
        minimum width=2.4cm, minimum height=2.5cm, text width=2.1cm,
        align=center, font=\tiny}
]
    % عنوان
    \node[draw=mahestan-purple, fill=bg-light-purple, rounded corners=6pt,
        minimum width=8cm, minimum height=0.8cm, align=center, font=\small\bfseries]
        (title) at (5.2,2.5) {⚖️ عدالت انتقالی ایران};

    % پنج ستون
    \node[pillar=mahestan-blue]   (p1) at (0,0)
        {\textbf{🔍 حقیقت‌یابی}\\[3pt] مستندسازی\\شهادت قربانیان\\جلسات علنی\\گزارش جامع};
    \node[pillar=mahestan-red]    (p2) at (2.7,0)
        {\textbf{⚔️ عدالت کیفری}\\[3pt] محاکمات عادلانه\\تفکیک سطوح\\مسئولیت\\عفو مشروط};
    \node[pillar=mahestan-orange] (p3) at (5.4,0)
        {\textbf{💰 جبران خسارت}\\[3pt] غرامت مالی\\اعاده‌ی حیثیت\\خدمات درمانی\\بازگرداندن اموال};
    \node[pillar=mahestan-gold]   (p4) at (8.1,0)
        {\textbf{🕊️ آشتی ملی}\\[3pt] گفتگوی ملی\\موزه‌ها و یادبودها\\اصلاح تاریخ‌نگاری\\روز ملی یادبود};
    \node[pillar=mahestan-green]  (p5) at (10.8,0)
        {\textbf{🛡️ عدم تکرار}\\[3pt] اصلاح نهادی\\آموزش حقوق بشر\\نهادهای نظارتی\\قانون اساسی};

    % فلش‌ها
    \foreach \p in {p1,p2,p3,p4,p5}
        \draw[-{Stealth}, thick, mahestan-purple!60] (title.south) -- (\p.north);
\end{tikzpicture}
\end{center}

% ==================== بخش ۵ ====================
\section{اولویت‌های ۱۰۰ روز نخست}

\subsection{هفته‌ی اول — ده فرمان اضطراری}

\begin{alertbox}[title={🔴 فرمان‌های اضطراری — بلافاصله پس از تشکیل مجلس}]
\begin{center}
\begin{tabularx}{\textwidth}{>{\centering\bfseries\small}p{0.5cm} >{\small}X >{\centering\small}p{1cm}}
\rowcolor{mahestan-red!15}
\textcolor{mahestan-red}{\#} & \textcolor{mahestan-red}{\bfseries اقدام} & \textcolor{mahestan-red}{\bfseries نماد} \\
\midrule
۱ & آزادی فوری تمام زندانیان سیاسی و عقیدتی & 🔓 \\
\rowcolor{bg-light-red}
۲ & لغو حجاب اجباری و کلیه‌ی قوانین کنترل پوشش & 👗 \\
۳ & رفع کامل فیلترینگ و سانسور اینترنت & 🌐 \\
\rowcolor{bg-light-red}
۴ & تعلیق مجازات اعدام & ⛔ \\
۵ & آزادی فعالیت احزاب، اتحادیه‌ها و تشکل‌ها & ✊ \\
\rowcolor{bg-light-red}
۶ & ممنوعیت بازداشت بدون حکم قضایی & ⚖️ \\
۷ & لغو محاکم انقلاب و دادگاه‌های ویژه‌ی روحانیت & 🏛️ \\
\rowcolor{bg-light-red}
۸ & اعلام حق بازگشت آزادانه‌ی تبعیدیان & ✈️ \\
۹ & مسدودسازی موقت دارایی‌های مشکوک نهادهای حکومتی & 🔒 \\
\rowcolor{bg-light-red}
۱۰ & حفاظت از اسناد و بایگانی‌های نهادهای امنیتی & 📁 \\
\end{tabularx}
\end{center}
\end{alertbox}

\subsection{هفته‌ی ۱ تا ۱۴ — سایر اولویت‌ها}

\begin{center}
\begin{tikzpicture}[
    phase/.style={draw=#1, fill=#1!10, rounded corners=5pt,
        minimum width=3.5cm, minimum height=3cm, text width=3.2cm,
        align=center, font=\tiny}
]
    \node[phase=mahestan-orange] (high) at (0,0)
        {\textbf{\small 🟠 هفته‌ی ۱–۴}\\[3pt]
        \textbf{اولویت بالا}\\[5pt]
        قانون آزادی مطبوعات\\
        قانون احزاب\\
        قانون تظاهرات\\
        قانون بحران اقتصادی\\
        شفافیت مالی مقامات\\
        لغو تبعیض علیه زنان\\
        حقوق اقلیت‌ها\\
        حفاظت داده‌ها};

    \node[phase=mahestan-gold] (med) at (4.5,0)
        {\textbf{\small 🟡 هفته‌ی ۴–۸}\\[3pt]
        \textbf{اولویت متوسط}\\[5pt]
        قانون عدالت انتقالی\\
        کمیسیون حقیقت‌یابی\\
        قانون اصلاح امنیتی\\
        استقلال قضایی\\
        استقلال بانک مرکزی\\
        حسابرسی ملی};

    \node[phase=mahestan-green] (norm) at (9,0)
        {\textbf{\small 🟢 هفته‌ی ۸–۱۴}\\[3pt]
        \textbf{اولویت عادی}\\[5pt]
        قانون محیط زیست\\
        اصلاح نظام آموزشی\\
        خصوصی‌سازی شفاف\\
        سرمایه‌گذاری خارجی\\
        تشکیل مجلس مؤسسان\\
        الحاق به کنوانسیون‌ها};

    \draw[-{Stealth}, very thick, mahestan-orange] (high) -- (med);
    \draw[-{Stealth}, very thick, mahestan-gold] (med) -- (norm);
\end{tikzpicture}
\end{center}

% ==================== بخش ۶ ====================
\section{اصول بنیادین غیرقابل‌نقض}

\begin{principlebox}[title={🛡️ ۱۴ اصل غیرقابل‌نقض — خط قرمز مطلق}]
\begin{multicols}{2}
\begin{enumerate}[label=\textcolor{mahestan-blue}{\bfseries\arabic*.}, itemsep=2pt, font=\small]
    \item 🗳️ \textbf{حاکمیت مردم}
    \item ⛪ \textbf{جدایی دین از دولت}
    \item ⚖️ \textbf{برابری همه در قانون}
    \item 🗣️ \textbf{آزادی‌های بنیادین}
    \item 🚫 \textbf{ممنوعیت مطلق شکنجه}
    \item 👨‍⚖️ \textbf{حق دادرسی عادلانه}
    \item 🏘️ \textbf{حقوق اقلیت‌ها}
    \item 👩 \textbf{حقوق و برابری زنان}
    \item ⚖️ \textbf{استقلال قضایی}
    \item 🗺️ \textbf{تمامیت ارضی + حق تعیین سرنوشت}
    \item 📜 \textbf{حاکمیت قانون}
    \item 🕊️ \textbf{عدالت نه انتقام}
    \item 🌐 \textbf{آزادی اطلاعات}
    \item ☮️ \textbf{تعهد به صلح}
\end{enumerate}
\end{multicols}
\vspace{0.2cm}
{\footnotesize\textcolor{text-medium}{%
\textbf{تضمین اجرا:}
دوران گذار: اراده‌ی مردم + نهادهای مدنی ناظر
\quad|\quad
نظام نهایی: دادگاه مستقل قانون اساسی}}
\end{principlebox}

% ==================== بخش ۷ ====================
\section{نقشه‌ی راه زمانی}

\begin{center}
\begin{tikzpicture}[
    timeline/.style={draw=#1, fill=#1!12, rounded corners=5pt,
        minimum width=1.8cm, minimum height=2.8cm, text width=1.6cm,
        align=center, font=\tiny},
    arr/.style={-{Stealth[length=4pt]}, very thick, #1}
]
    \node[timeline=border-gray] (p0) at (0,0)
        {\textbf{\footnotesize فاز ۰}\\[2pt]
        \textcolor{text-medium}{پیش از گذار}\\[3pt]
        پیش‌نویس‌ها\\ائتلاف‌سازی\\طراحی زیرساخت};
    \node[timeline=mahestan-orange] (p1) at (2.3,0)
        {\textbf{\footnotesize فاز ۱}\\[2pt]
        \textcolor{mahestan-orange}{روز ۰ → هفته ۵}\\[3pt]
        انتخابات\\تشکیل مجلس\\مهستان};
    \node[timeline=mahestan-red] (p2) at (4.6,0)
        {\textbf{\footnotesize فاز ۲}\\[2pt]
        \textcolor{mahestan-red}{هفته ۱–۸}\\[3pt]
        فرمان‌های\\اضطراری\\کابینه};
    \node[timeline=mahestan-purple] (p3) at (6.9,0)
        {\textbf{\footnotesize فاز ۳}\\[2pt]
        \textcolor{mahestan-purple}{هفته ۱–۱۰}\\[3pt]
        ق.ا. موقت\\+ رفراندوم};
    \node[timeline=mahestan-blue] (p4) at (9.2,0)
        {\textbf{\footnotesize فاز ۴-۵}\\[2pt]
        \textcolor{mahestan-blue}{ماه ۳–۱۸}\\[3pt]
        نهادسازی\\ق.ا. نهایی\\+ رفراندوم};
    \node[timeline=mahestan-green] (p5) at (11.5,0)
        {\textbf{\footnotesize فاز ۶}\\[2pt]
        \textcolor{mahestan-green}{ماه ۱۸–۲۴}\\[3pt]
        انتخابات نهایی\\انتقال قدرت\\پایان گذار};

    \draw[arr=border-gray] (p0) -- (p1);
    \draw[arr=mahestan-orange] (p1) -- (p2);
    \draw[arr=mahestan-red] (p2) -- (p3);
    \draw[arr=mahestan-purple] (p3) -- (p4);
    \draw[arr=mahestan-blue] (p4) -- (p5);
\end{tikzpicture}
\end{center}

{\footnotesize\textcolor{text-medium}{%
⚠️ دوره‌ی مجلس: حداکثر ۲ سال + یک تمدید ۶ ماهه با رأی ⅔ نمایندگان.
تمام زمان‌بندی‌ها نسبی و از روز صفر (آغاز رسمی گذار) محاسبه می‌شوند.}}

% ==================== بخش ۸ ====================
\section{فراخوان مشارکت}

\begin{greenbox}[title={📣 دعوت از همه‌ی ایرانیان آزادیخواه}]
\begin{itemize}[itemsep=6pt, font=\small]
    \item[\textcolor{mahestan-blue}{🏛️}] \textbf{احزاب و جریان‌ها:} از هم‌اکنون پیش‌نویس قانون اساسی موقت پیشنهادی‌تان را تدوین و منتشر کنید.
    \item[\textcolor{mahestan-green}{👤}] \textbf{هر شهروند:} مطالعه کنید، بحث کنید، نقد سازنده کنید. مشارکت کنید.
    \item[\textcolor{mahestan-orange}{🎓}] \textbf{هر متخصص:} دانش و تجربه‌تان را در خدمت طراحی نهادی گذار بگذارید.
    \item[\textcolor{mahestan-red}{✊}] \textbf{هر فعال حقوق بشر:} مستندسازی را ادامه دهید — اسناد امروز، مبنای عدالت فرداست.
\end{itemize}
\end{greenbox}

\vspace{0.5cm}

% نمودار پایانی
\begin{center}
\begin{tikzpicture}
    \node[draw=mahestan-red, fill=mahestan-red!12, rounded corners=6pt,
        minimum width=3cm, minimum height=1.5cm, align=center, font=\small]
        (past) at (0,0) {\textbf{\textcolor{mahestan-red}{🔴 گذشته}}\\[2pt]
        {\tiny استبداد · سرکوب\\تبعیض · فساد · انزوا}};

    \node[draw=mahestan-gold, fill=mahestan-gold!15, rounded corners=6pt,
        minimum width=3.5cm, minimum height=1.5cm, align=center, font=\small]
        (trans) at (5,0) {\textbf{\textcolor{mahestan-gold}{🟡 گذار}}\\[2pt]
        {\tiny مجلس مهستان\\🔬 تخصص + 📢 دموکراسی\\عدالت انتقالی · مشارکت}};

    \node[draw=mahestan-green, fill=mahestan-green!12, rounded corners=6pt,
        minimum width=3cm, minimum height=1.5cm, align=center, font=\small]
        (future) at (10.5,0) {\textbf{\textcolor{mahestan-green}{🟢 آینده}}\\[2pt]
        {\tiny دموکراسی · حقوق بشر\\برابری · رفاه · صلح}};

    \draw[-{Stealth[length=6pt]}, ultra thick, mahestan-gold] (past) -- (trans);
    \draw[-{Stealth[length=6pt]}, ultra thick, mahestan-green] (trans) -- (future);
\end{tikzpicture}
\end{center}

\vspace{0.5cm}

\begin{center}
    \textcolor{mahestan-purple}{\rule{6cm}{1.5pt}}\\[0.4cm]
    {\Large\bfseries\textcolor{mahestan-blue}{«ایران برای همه‌ی ایرانیان»}}\\[0.4cm]
    \textcolor{mahestan-purple}{\rule{6cm}{1.5pt}}\\[0.4cm]
    {\small جزئیات بیشتر در کتاب \textbf{«شکوفایی: ایران از بحران تا بالندگی»}}\\[0.2cm]
    {\footnotesize\textcolor{text-medium}{این سند زنده است و به‌روزرسانی خواهد شد. آخرین ویرایش: اسفند ۱۴۰۴}}
\end{center}

\end{document}